% !TEX program = xelatex
% 数学物理方法1 · 第四章 · 作业 3 详解与约当引理的应用（图片版）
\documentclass[11pt]{ctexart}
%=====================================================================
% theme-violet.tex —— 讲义精讲 · 紫花主题
% 用法：\documentclass[11pt]{ctexart}
%       \input{assets/theme-violet}
%       \begin{document} ... \end{document}
% 只用 ctexart 自带能力 + xcolor/enumitem/fancyhdr/pifont/amsmath，
% 不依赖 tikz/tcolorbox，保证 MiKTeX 可直接编译。
%=====================================================================

\usepackage[a4paper,left=2.0cm,right=2.0cm,top=2.1cm,bottom=2.3cm,%
            headsep=12pt,footskip=22pt]{geometry}
\usepackage{amsmath,amssymb}
\usepackage{needspace}
\usepackage{xcolor}
\usepackage{array}
\usepackage{enumitem}
\usepackage{pifont}
\usepackage{fancyhdr}
\usepackage[hidelinks]{hyperref}

%---------------------------------------------------------------- 配色
\definecolor{violet}{RGB}{107,78,143}      % 主色：紫
\definecolor{violetdeep}{RGB}{63,44,92}    % 深紫：标题/强调
\definecolor{petal}{RGB}{183,159,208}      % 花瓣：装饰线、花饰
\definecolor{bandviolet}{RGB}{247,243,251} % 浅紫底：档案/提示框
\definecolor{bandcream}{RGB}{252,249,242}  % 米白底：易错/注意框
\definecolor{rulepurple}{RGB}{217,207,230} % 分隔线
\definecolor{inkgray}{RGB}{92,96,110}      % 次要文字

%---------------------------------------------------------------- 字体
% 有 Noto CJK 用 Noto，没有就退回 ctex 默认（Windows 下为宋体/黑体）
\IfFontExistsTF{Noto Serif CJK SC}{\setCJKmainfont{Noto Serif CJK SC}}{}
\IfFontExistsTF{Noto Sans CJK SC}{\setCJKsansfont{Noto Sans CJK SC}}{}

%---------------------------------------------------------------- 版式
\setlist{nosep,leftmargin=1.9em}
\linespread{1.28}
\setlength{\parskip}{3.5pt}
\setlength{\parindent}{0pt}
\setlength{\tabcolsep}{5pt}
\renewcommand{\arraystretch}{1.28}

%---------------------------------------------------------------- 花饰
% \ding{95} = 五瓣小花，\ding{96} = 星形花；都来自 pifont，必定存在
\newcommand{\orn}{\textcolor{petal}{\ding{95}}}          % 行内小标
\newcommand{\ornstar}{\textcolor{petal}{\ding{96}}}      % 星形花
\newcommand{\flowerrule}{%                                % 居中花饰分隔线
  \par\vspace{7pt}%
  \noindent\makebox[\linewidth]{\color{petal}%
    \rule{0.20\linewidth}{0.6pt}\hspace{8pt}{\small\ding{95}}\hspace{4pt}%
    \raisebox{1.1pt}{\tiny\ding{95}}\hspace{4pt}{\small\ding{95}}\hspace{8pt}%
    \rule{0.20\linewidth}{0.6pt}}%
  \par\vspace{7pt}}

%---------------------------------------------------------------- 标题与节头
\newcommand{\doctitle}[3]{%   #1 主标题  #2 副标题  #3 元信息（可用 \\ 分行）
  \begin{center}
    {\color{petal}\small \orn\enspace #2\enspace \ornstar}\\[7pt]
    {\LARGE\bfseries\color{violetdeep} #1}\\[10pt]
    {\color{petal}\rule{0.55\linewidth}{1pt}}\\[4pt]
    {\small\color{inkgray} #3}
  \end{center}\vspace{2pt}}

\newcommand{\hdr}[2]{%   #1 章节号（如 §0）  #2 标题
  \par\vspace{12pt}%
  \noindent\colorbox{violetdeep}{\color{white}\bfseries\small\;#1\;}\hspace{7pt}%
  {\color{violetdeep}\large\bfseries #2}\par\nobreak
  \vspace{3pt}{\color{petal}\hrule height 1.2pt}\vspace{6pt}}

\newcommand{\subhdr}[1]{%
  \par\vspace{6pt}%
  \noindent{\color{violet}\bfseries #1}\par\vspace{1pt}}

%---------------------------------------------------------------- 提示框
\newcommand{\callout}[1]{% 浅紫底：讲义档案、读法说明
  \par\medskip\noindent
  \fcolorbox{rulepurple}{bandviolet}{\parbox{\dimexpr\linewidth-2\fboxsep-2\fboxrule}{\vspace{2pt}#1\vspace{2pt}}}%
  \par\medskip}

\newcommand{\warnbox}[1]{% 米白底：易错、边界条件
  \par\medskip\noindent
  \fcolorbox{rulepurple}{bandcream}{\parbox{\dimexpr\linewidth-2\fboxsep-2\fboxrule}{\vspace{2pt}#1\vspace{2pt}}}%
  \par\medskip}

\newcommand{\gapitem}[2]{% #1 缺口类型（紫） #2 内容
  \par\vspace{2pt}\noindent
  \hangindent=1.4em\hangafter=1
  {\color{petal}\small\ding{96}}\enspace{\bfseries\color{violetdeep}#1}\enspace #2\par}

% 题干块（\qstem）：每道例题／作业的「详解」开头必须先复述题干与出处。
% 起因：2026-10-25 第五章讲义把七道作业的题面集中放在一页、详解另起一页，
% 读者翻到「第 4 题」时看不到条件，等于一道没有题干的题目。
% 规则：详解块内自带题干，读者不翻页就能读懂题在问什么。
\newcommand{\qstem}[1]{%
  \par\medskip\noindent
  \fcolorbox{violet}{bandviolet}{%
    \parbox{\dimexpr\linewidth-2\fboxsep-2\fboxrule}{%
      \vspace{2pt}%
      {\small\bfseries\color{violet}题干}\hspace{0.8em}#1%
      \vspace{2pt}}}%
  \par\medskip}

% v2: 保护题干与开头，长题可跨页，续页补上下文。
\newcommand{\qbreak}{\needspace{6\baselineskip}}
% 独立练习：ID、来源、题干、考点、答案标签（通常 ans:ID）。
\newcommand{\practice}[5]{%
  \qbreak\subhdr{练习 #1}\qstem{（#2）#3}%
  \textbf{考点}：#4\par
  详解见附录第~\pageref{#5}~页。\par\medskip}
% 附录开头：ID、来源、条件摘要；自动产生 ans:ID 标签。
\newcommand{\answerstart}[3]{%
  \qbreak\subhdr{#1\quad 详解}\label{ans:#1}\qstem{（#2）#3}}
\newcommand{\answercontinue}[2]{%
  \needspace{5\baselineskip}\subhdr{#1\quad 续解}%
  {\small\color{inkgray}本页接续：#2}\par\medskip}
% 延迟提供默认值，让旧讲义可在导言区保留自己的 \newcommand{\step}。
\AtBeginDocument{\providecommand{\step}[1]{\needspace{3\baselineskip}\par\medskip\textbf{#1}\par}}

% 前提旁注：只放在「需要理解的那一步」后面，给陈述 + 本题代入，不重证。
\newcommand{\aside}[1]{%
  \par\vspace{1pt}\noindent
  \hangindent=1.6em\hangafter=1
  {\small\color{inkgray}%
    \hspace*{0.6em}{\color{petal}$\triangleright$}\hspace{2pt}#1}%
  \par\vspace{2pt}}

%---------------------------------------------------------------- 表格
% 本机 MiKTeX 的 colortbl 与 array 版本不匹配（\rowcolor 会报错），
% 故表格统一用「上粗紫线 + 表头 + 中细线」的三线式，不依赖 colortbl。
\newcommand{\tabtop}{\par\vspace{3pt}{\color{violet}\hrule height 1.2pt}\vspace{1pt}}
\newcommand{\tabmid}{\par\vspace{1pt}{\color{petal}\hrule height 0.5pt}\vspace{1pt}}
\newcommand{\tabbot}{\par\vspace{3pt}{\color{rulepurple}\hrule height 0.5pt}\vspace{8pt}}
\newcommand{\thead}[1]{{\bfseries\color{violetdeep}#1}}

%---------------------------------------------------------------- 页眉页脚
\pagestyle{fancy}
\fancyhf{}
\renewcommand{\headrulewidth}{0.4pt}
\renewcommand{\headrule}{\hbox to\headwidth{\color{rulepurple}\leaders\hrule height \headrulewidth\hfill}}
\fancyhead[L]{\small\color{inkgray}\songti \doctag}
\fancyhead[R]{\small\color{inkgray}\songti \coursetag}
\fancyfoot[C]{\small\color{inkgray}\orn\enspace{\color{violet}\thepage}\enspace\orn}

% 每篇讲义在导言区设这两个值
\newcommand{\doctag}{}
\newcommand{\coursetag}{}
\newcommand{\settags}[2]{\renewcommand{\doctag}{#1}\renewcommand{\coursetag}{#2}}

%---------------------------------------------------------------- 数学记号
\newcommand{\vect}[1]{\boldsymbol{#1}}
\newcommand{\dd}{\,\mathrm{d}}
\newcommand{\ddx}{\frac{\mathrm{d}}{\mathrm{d}x}}
\addtolength{\jot}{2pt}



\settags{第四章 · 作业 3 详解}{数学物理方法1}

\renewcommand{\qbreak}{\needspace{0.45\textheight}}
\setlength{\headheight}{22pt}
\setlength{\abovedisplayskip}{6pt plus 2pt minus 2pt}
\setlength{\belowdisplayskip}{6pt plus 2pt minus 2pt}
\setlength{\abovedisplayshortskip}{3pt}
\setlength{\belowdisplayshortskip}{3pt}

\begin{document}

\doctitle{第四章　作业 3 详解与约当引理}
{三道有理积分，一条补弧，一台“约当机器”}
{依据课件 p62（作业 3）与 p53--p54（约当引理）\quad 2026-11-03\quad 数学物理方法1}

\callout{\textbf{先说结论（也是这份材料的主线）}：作业 3 的三道题都是\textbf{有理式}的无穷积分，被积函数里没有 $e^{imx}$，所以它们真正满足的是\textbf{大圆弧引理}（判据 $|zf(z)|\to0$）。约当引理是它的“带振荡因子”加强版，两者\textbf{共用同一套证明骨架}——参数化、取模、把圆弧上的积分压到零。本材料把这套骨架完整拆开（§1），说明它在 $m=0$ 时如何\textbf{退化成}大圆弧引理（§2），再用它逐题算完（§3--§5），最后给一道\textbf{非约当引理不可}的对照题（§6）。}

\hdr{§0}{三题答案与“该用哪条引理”}

\begin{center}
\small
\tabtop
\begin{tabular}{@{}cll@{}}
\thead{题} & \thead{积分} & \thead{结果} \\
\end{tabular}
\tabmid
\begin{tabular}{@{}cll@{}}
(1) & $\displaystyle\int_{-\infty}^{\infty}\frac{x^{2}\dd x}{(x^{2}+1)(x^{2}+4)}$ & $\dfrac{\pi}{3}$ \\
(2) & $\displaystyle\int_{0}^{\infty}\frac{x^{2}\dd x}{x^{4}+a^{4}}\ (a>0)$ & $\dfrac{\pi}{2\sqrt2\,a}=\dfrac{\pi\sqrt2}{4a}$ \\
(3) & $\displaystyle\int_{-\infty}^{\infty}\frac{x^{2}\dd x}{(x^{2}+a^{2})^{2}}\ (a>0)$ & $\dfrac{\pi}{2a}$ \\
\end{tabular}
\tabbot
\end{center}

三题是同一个流程：把实轴段 $[-R,R]$ 补上半圆弧 $C_R$ 围成闭围道，说明\textbf{圆弧上的积分在 $R\to\infty$ 时归零}（引理就守在这一步），再用留数定理，最后取极限。

\subhdr{0.1\quad 三条引理的分工}

\begin{center}
\small
\tabtop
\begin{tabular}{@{}lp{0.30\linewidth}p{0.34\linewidth}@{}}
\thead{被积函数} & \thead{判据} & \thead{用在哪一步} \\
\end{tabular}
\tabmid
\begin{tabular}{@{}lp{0.30\linewidth}p{0.34\linewidth}@{}}
有理式 $\dfrac{P}{Q}$，$\deg Q-\deg P\ge2$ & 大圆弧引理：$R\max\limits_{C_R}|f|\to0$（等价 $|zf|\to0$） & 弧上积分归零 \\
带 $e^{imx},\cos mx,\sin mx$（$m>0$） & 约当引理：$\max\limits_{C_R}|f|\to0$ 一致，$m>0$ & 弧上积分归零 \\
实轴上有单极点 & 小圆弧引理 & 得 $\pm i\pi\operatorname{Res}$ \\
\end{tabular}
\tabbot
\end{center}

\begin{center}
\fcolorbox{violet}{bandviolet}{\parbox{0.92\linewidth}{\centering\vspace{3pt}
\textbf{\color{violetdeep}约当引理 $=$ 大圆弧引理 $+$ 一个振荡因子}\\[3pt]
\small 有振荡因子：$|f|$ 只要一致趋于 $0$ 就够（哪怕只按 $1/R$ 掉）。\\
\small 无振荡因子：必须 $R\max|f|\to0$，即分母比分子至少高\textbf{两次}。\vspace{3pt}}}
\end{center}

所以卷面上写“由大圆弧引理（$|zf(z)|\to0$）”最稳；若写“由约当引理”，必须补一句 $\max|f|\to0$ 一致的验证，并说明取的是振荡因子为 $1$ 的退化情形。

\hdr{§1}{约当引理：陈述、证明、四条边界}

\subhdr{1.1\quad 陈述（课件 p53）}

设 $m>0$，$f$ 在 $C_R$ 上连续，且
\[
\max_{z\in C_R}|f(z)|\longrightarrow0\qquad(R\to\infty),
\]
则
\[
\lim_{R\to\infty}\int_{C_R}f(z)e^{imz}\dd z=0 .
\]
课件写“在上半平面及实轴上一致趋于零”，其实\textbf{只要弧上一致}就够，因为估计只用到弧上的值。

\subhdr{1.2\quad 证明的四个动作（可直接抄进卷面）}

\textbf{动作一：参数化。}在 $C_R$ 上 $z=Re^{i\theta}$，$\theta:0\to\pi$，$\dd z=iRe^{i\theta}\dd\theta$，故 $|\dd z|=R\dd\theta$。

\textbf{动作二：取模，并把 $e^{imz}$ 按实部、虚部拆开。}
\[
\bigl|e^{imz}\bigr|=\bigl|e^{im(R\cos\theta+iR\sin\theta)}\bigr|
=\underbrace{\bigl|e^{imR\cos\theta}\bigr|}_{=\,1}\cdot
\underbrace{\bigl|e^{-mR\sin\theta}\bigr|}_{=\,e^{-mR\sin\theta}}
=e^{-mR\sin\theta}.
\]
实部只让因子在单位圆上转，模恒为 $1$；真正起作用的是\textbf{虚部}，它给出指数衰减。又因 $\theta\in[0,\pi]$ 时 $\sin\theta\ge0$，所以 $e^{-mR\sin\theta}\le1$。于是
\[
\left|\int_{C_R}fe^{imz}\dd z\right|
\le\int_0^{\pi}|f(Re^{i\theta})|\,e^{-mR\sin\theta}R\dd\theta
\le M_RR\int_0^{\pi}e^{-mR\sin\theta}\dd\theta ,
\]
其中 $M_R=\max_{C_R}|f|$。

\textbf{动作三：用 $\sin\theta\ge\dfrac{2\theta}{\pi}$ 把指数化成线性，就能积出来。}
\[
\int_0^{\pi}e^{-mR\sin\theta}\dd\theta
=2\int_0^{\pi/2}e^{-mR\sin\theta}\dd\theta
\le2\int_0^{\pi/2}e^{-2mR\theta/\pi}\dd\theta
=2\cdot\frac{\pi}{2mR}\bigl(1-e^{-mR}\bigr).
\]

\textbf{动作四：取 $R\to\infty$。}代入得
\[
\left|\int_{C_R}fe^{imz}\dd z\right|
\le M_RR\cdot\frac{\pi}{mR}\bigl(1-e^{-mR}\bigr)
=\frac{\pi M_R}{m}\bigl(1-e^{-mR}\bigr)\le\frac{\pi M_R}{m}\longrightarrow0 .
\]
注意 $R$ \textbf{被约掉了}；约掉它靠的是 $m>0$，因为只有 $m>0$ 时 $e^{-mR}$ 才随 $R$ 掉下去。\hfill$\square$

\subhdr{1.3\quad 那个不等式从哪来}

$\sin\theta$ 在 $[0,\pi/2]$ 上上凸（$\sin''\theta=-\sin\theta<0$），图像在弦的上方；连接 $(0,0)$ 与 $(\pi/2,1)$ 的弦正是 $y=2\theta/\pi$，故 $\sin\theta\ge2\theta/\pi$。等价地，$\dfrac{\sin\theta}{\theta}$ 在该区间递减，右端点取到最小值 $2/\pi$。这个常数就是动作三里 $2mR/\pi$ 的来历。

\warnbox{\textbf{四条边界（扣分高发）}\\[2pt]
1. \textbf{$m<0$ 不能走上半平面}：$|e^{imz}|=e^{-mR\sin\theta}=e^{|m|R\sin\theta}$ 在上半平面爆炸。应改走下半个半圆（那里 $\sin\theta<0$），方向变顺时针，\textbf{结果多一个负号}。\\
2. \textbf{$|f|$ 不一致趋于零则失效}：振荡只能压住不衰减的模，压不住发散的模。\\
3. \textbf{实轴上有极点不能直接闭合}：要先挖小半圆，此时围道极限是 $2\pi i\sum_{\text{上半}}\operatorname{Res}\pm i\pi\sum_{\text{实轴}}\operatorname{Res}$。\\
4. \textbf{不能“令 $m=0$”反推}：$m=0$ 时 $e^{imz}\equiv1$，动作二、三整条链断掉，只剩 ML 估计——这就是下一节的退化情形。}

\hdr{§2}{同一台机器在 $m=0$：退化为大圆弧引理}

把 §1 的动作一、二照做，但被积函数是 $f$ 而不是 $fe^{imz}$（等价于 $m=0$、$|e^{imz}|=1$）。这时振荡压缩全没了，只剩 ML 估计：
\[
\left|\int_{C_R}f(z)\dd z\right|\le M_R\cdot\pi R ,
\qquad\text{判据：}\quad R\max_{C_R}|f(z)|\longrightarrow0 .
\]
这正是课件 p47 的 $|zf(z)|\to0$（令 $|z|=R$ 就是同一件事）。对有理式 $f=P/Q$，$\deg Q=n$、$\deg P=m$ 时 $|f|\sim C/R^{\,n-m}$，故
\[
RM_R\sim C\,R^{\,1-(n-m)}\longrightarrow0
\iff n-m\ge2 .
\]
\textbf{这就是“次数差至少两次”的定量来源}，也是三题都能补弧的理由：三题的次数差都恰好是 $4-2=2$（临界值），所以 $RM_R\sim\pi/R$，收敛得慢，但够用。

\hdr{§3}{题 (1)：$\displaystyle I_1=\int_{-\infty}^{\infty}\frac{x^{2}\dd x}{(x^{2}+1)(x^{2}+4)}$}

\qstem{计算 $\displaystyle I_1=\int_{-\infty}^{\infty}\frac{x^{2}\dd x}{(x^{2}+1)(x^{2}+4)}$。}

\textbf{① 解析性。}$f(z)=\dfrac{z^{2}}{(z-i)(z+i)(z-2i)(z+2i)}$，四个一阶极点 $\pm i,\pm2i$ 全在虚轴上，\textbf{实轴无奇点}，此外处处解析。

\textbf{② 补弧。}$\Gamma_R=[-R,R]\cup C_R$（$R>2$），正方向：先沿实轴 $-R\to R$，再沿上半圆弧 $R\to-R$。留数定理给出
\[
\oint_{\Gamma_R}f\dd z=\int_{-R}^{R}f(x)\dd x+\int_{C_R}f\dd z
=2\pi i\!\!\sum_{z_k\in\text{上半平面}}\!\!\operatorname{Res}f(z_k).
\]

\textbf{③ 弧上归零（引理在本题的全部内容）。}在 $C_R$ 上 $|\dd z|=R\dd\theta$，且
\[
|z^{2}+1|\ge R^{2}-1,\quad|z^{2}+4|\ge R^{2}-4
\Longrightarrow
|f|\le\frac{R^{2}}{(R^{2}-1)(R^{2}-4)}=M_R ,
\]
\[
\left|\int_{C_R}f\dd z\right|\le\pi RM_R=\frac{\pi R^{3}}{(R^{2}-1)(R^{2}-4)}\sim\frac{\pi}{R}\longrightarrow0 .
\]
判据 $RM_R\to0$ 成立，即 $|zf(z)|\to0$：\textbf{大圆弧引理}。若按“约当引理”的写法讲，这里取的是 $m=0$ 的退化形式，动作二里 $|e^{imz}|\le1$ 这一步退化成纯 ML 估计，而 $M_R\to0$ 一致正是必须补写的那句验证。

\textbf{④ 留数。}上半平面一阶极点 $z=i,\;2i$：
\[
\operatorname{Res}_{z=i}f=\frac{i^{2}}{(2i)(i^{2}+4)}=\frac{-1}{6i}=\frac{i}{6},
\qquad
\operatorname{Res}_{z=2i}f=\frac{(2i)^{2}}{\bigl((2i)^{2}+1\bigr)(4i)}=\frac{-4}{-12i}=-\frac{i}{3}.
\]

\textbf{⑤ 取极限。}留数和 $=\dfrac{i}{6}-\dfrac{i}{3}=-\dfrac{i}{6}$，故
\[
I_1=2\pi i\left(-\frac{i}{6}\right)=\frac{\pi}{3}.
\]

\textbf{⑥ 校验（部分分式）。}$\dfrac{x^{2}}{(x^{2}+1)(x^{2}+4)}=-\dfrac{1/3}{x^{2}+1}+\dfrac{4/3}{x^{2}+4}$，于是
$I_1=-\dfrac{\pi}{3}+\dfrac{4}{3}\cdot\dfrac{\pi}{2}=\dfrac{\pi}{3}$。\hfill\checkmark

\hdr{§4}{题 (2)：$\displaystyle I_2=\int_{0}^{\infty}\frac{x^{2}\dd x}{x^{4}+a^{4}}\ (a>0)$}

\qstem{计算 $\displaystyle I_2=\int_{0}^{\infty}\frac{x^{2}\dd x}{x^{4}+a^{4}}$，$a>0$。}

\textbf{① 解析性。}极点由 $z^{4}=-a^{4}=a^{4}e^{i\pi}$ 给出：$z_k=ae^{i(2k+1)\pi/4}$，$k=0,1,2,3$，与实轴成 $45^\circ$，\textbf{实轴无奇点}。$a>0$ 保证上半平面只有 $z_1=ae^{i\pi/4}$、$z_2=ae^{3i\pi/4}$。

\textbf{② 偶性折半 + 补弧。}被积函数是偶函数，故 $I_2=\dfrac12\displaystyle\int_{-\infty}^{\infty}\dfrac{x^{2}\dd x}{x^{4}+a^{4}}$，再对右端补上半大圆弧（$R>a$）。

\textbf{③ 弧上归零。}$|z^{4}+a^{4}|\ge R^{4}-a^{4}$，故
\[
|f|\le\frac{R^{2}}{R^{4}-a^{4}}=M_R,
\qquad
\left|\int_{C_R}f\dd z\right|\le\pi RM_R=\frac{\pi R^{3}}{R^{4}-a^{4}}\sim\frac{\pi}{R}\longrightarrow0 .
\]
与题 (1) 完全同构：次数差 $4-2=2$，取到临界值。

\textbf{④ 留数（$\frac{1}{4z_k}$ 技巧）。}一阶极点处
$\operatorname{Res}_{z_k}\dfrac{z^{2}}{z^{4}+a^{4}}=\dfrac{z_k^{2}}{4z_k^{3}}=\dfrac{1}{4z_k}$，于是
\[
\sum_{k=1,2}\operatorname{Res}=\frac{1}{4a}\left(e^{-i\pi/4}+e^{-3i\pi/4}\right)
=\frac{1}{4a}\left(\frac{1-i}{\sqrt2}+\frac{-1-i}{\sqrt2}\right)=-\frac{i\sqrt2}{4a}.
\]

\textbf{⑤ 取极限。}
\[
\int_{-\infty}^{\infty}\frac{x^{2}\dd x}{x^{4}+a^{4}}=2\pi i\left(-\frac{i\sqrt2}{4a}\right)=\frac{\pi\sqrt2}{2a},
\qquad
I_2=\frac12\cdot\frac{\pi\sqrt2}{2a}=\frac{\pi}{2\sqrt2\,a}.
\]

\textbf{⑥ 校验。}被积函数量纲为 $1/$长度，结果必形如 $C/a$，与公式一致；取 $a=2$，公式给 $\sqrt2\pi/8\approx0.55536$，数值积分相同。\hfill\checkmark

\hdr{§5}{题 (3)：$\displaystyle I_3=\int_{-\infty}^{\infty}\frac{x^{2}\dd x}{(x^{2}+a^{2})^{2}}\ (a>0)$}

\qstem{计算 $\displaystyle I_3=\int_{-\infty}^{\infty}\frac{x^{2}\dd x}{(x^{2}+a^{2})^{2}}$，$a>0$。}

\textbf{① 解析性。}$f(z)=\dfrac{z^{2}}{(z-ia)^{2}(z+ia)^{2}}$，$z=\pm ia$ 都是\textbf{二阶极点}，实轴无奇点；上半平面只有 $z=ia$。

\textbf{② 补弧。}同前，$\Gamma_R=[-R,R]\cup C_R$（$R>a$）。

\textbf{③ 弧上归零。}$|z^{2}+a^{2}|\ge R^{2}-a^{2}$，故
\[
|f|\le\frac{R^{2}}{(R^{2}-a^{2})^{2}}=M_R,
\qquad
\left|\int_{C_R}f\dd z\right|\le\pi RM_R=\frac{\pi R^{3}}{(R^{2}-a^{2})^{2}}\sim\frac{\pi}{R}\longrightarrow0 .
\]
\textbf{这里有个陷阱}：分母带平方，看着“掉得更快”，其实次数仍是 $4$，$\pi R^{3}/R^{4}=1/R$，与题 (1)(2) 一模一样——临界，但成立。

\textbf{④ 留数（二阶极点规则）。}
\[
\operatorname{Res}_{z=ia}f=\lim_{z\to ia}\frac{\dd}{\dd z}\left[\frac{z^{2}}{(z+ia)^{2}}\right]
=\lim_{z\to ia}\frac{2z(z+ia)-2z^{2}}{(z+ia)^{3}} .
\]
在 $z=ia$ 处：分子 $=2(ia)(2ia)-2(ia)^{2}=-4a^{2}+2a^{2}=-2a^{2}$；分母 $=(2ia)^{3}=-8ia^{3}$。故
\[
\operatorname{Res}_{z=ia}f=\frac{-2a^{2}}{-8ia^{3}}=\frac{1}{4ia}=-\frac{i}{4a},
\qquad
I_3=2\pi i\left(-\frac{i}{4a}\right)=\frac{\pi}{2a}.
\]

\textbf{⑤ 校验（代数拆分）。}用 $x^{2}=(x^{2}+a^{2})-a^{2}$：
\[
I_3=\int_{-\infty}^{\infty}\frac{\dd x}{x^{2}+a^{2}}-a^{2}\int_{-\infty}^{\infty}\frac{\dd x}{(x^{2}+a^{2})^{2}}
=\frac{\pi}{a}-a^{2}\cdot\frac{\pi}{2a^{3}}=\frac{\pi}{2a} .\quad\checkmark
\]
顺手记住这两个常用值：$\displaystyle\int_{-\infty}^{\infty}\frac{\dd x}{x^{2}+a^{2}}=\frac{\pi}{a}$，$\displaystyle\int_{-\infty}^{\infty}\frac{\dd x}{(x^{2}+a^{2})^{2}}=\frac{\pi}{2a^{3}}$。

\hdr{§6}{对照：什么时候非约当引理不可}

同一条围道、同一套留数，只把被积函数换掉，就能看出两条引理的边界。

\textbf{对照 A（仍是大圆弧）。}$\displaystyle\int_{-\infty}^{\infty}\frac{x\sin mx}{(x^{2}+1)(x^{2}+4)}\dd x$。取 $g=\dfrac{ze^{imz}}{(x^{2}+1)(x^{2}+4)}$ 型的复函数，在 $C_R$ 上 $|g|\sim1/R^{3}$，$\pi R\max|g|\sim\pi/R^{2}\to0$，\textbf{大圆弧引理就够}，$e^{imz}$ 只在最后取虚部时用。

\textbf{对照 B（非约当不可）。}$\displaystyle\int_{-\infty}^{\infty}\frac{x\sin mx}{x^{2}+1}\dd x$。取 $g(z)=\dfrac{ze^{imz}}{z^{2}+1}$，在 $C_R$ 上
\[
|g|\le\frac{R}{R^{2}-1}\sim\frac1R,
\qquad
R\max|g|\sim\frac{R^{2}}{R^{2}-1}\longrightarrow1\neq0 .
\]
大圆弧引理\textbf{失效}（判据要 $RM_R\to0$，这里趋于 $1$），弧上积分确实不趋于零。此时只有约当引理能救：把振荡因子单独利用，
\[
\left|\int_{C_R}ge^{imz}\dd z\right|
\le\int_0^\pi\frac{R}{R^{2}-1}e^{-mR\sin\theta}R\dd\theta
\le\frac{R^{2}}{R^{2}-1}\cdot\frac{\pi}{mR}\bigl(1-e^{-mR}\bigr)
\le\frac{\pi}{m}\cdot\frac{R}{R^{2}-1}\longrightarrow0 .
\]
顺序与 §1 完全一样，$R$ 被约掉靠的是 $m>0$。留数 $\operatorname{Res}_{z=i}g=\dfrac{ie^{-m}}{2i}=\dfrac{e^{-m}}{2}$，故
\[
\int_{-\infty}^{\infty}\frac{x\sin mx}{x^{2}+1}\dd x=\pi e^{-m}
\qquad(m=2\ \text{时}=0.425168,\ \text{与数值积分一致}).
\]

\flowerrule

\hdr{§7}{考场模板与三个坑}

\textbf{每题五句话：}\;(1)“$f$ 在实轴无奇点，上半平面有极点 $z_1,\dots,z_n$，此外解析。”\;(2)“取 $\Gamma_R=[-R,R]\cup C_R$，由留数定理 $\oint_{\Gamma_R}f\dd z=2\pi i\sum\operatorname{Res}$。”\;(3)“在 $C_R$ 上 $|f|\le M_R=\cdots$，故 $|\int_{C_R}f\dd z|\le\pi RM_R=\cdots\to0$（大圆弧引理／约当引理）。”\;(4) 逐个算留数，注明阶数与所用公式。\;(5)“令 $R\to\infty$，得 $I=2\pi i\sum\operatorname{Res}=\cdots$。”

\warnbox{\textbf{三个扣分点}\\[2pt]
1. 只说“$R$ 很大所以弧上积分为零”，不给 $M_R$ 与 $\pi RM_R$ 的显式估计——判据必须写出来。\\
2. 二阶极点误用 $\dfrac{P(z_0)}{Q'(z_0)}$：那条只对一阶极点成立；二阶要用 $\lim\dfrac{\dd}{\dd z}\bigl[(z-z_0)^{2}f\bigr]$，题 (3) 正是这一条。\\
3. 偶函数折半忘了说：题 (2) 必须写“被积函数是偶函数，故 $=\frac12\int_{-\infty}^{\infty}$”再补弧，否则围道与积分区间对不上。}

\vspace{4pt}
{\small\color{inkgray} 三题结果经独立复核：$\pi/3$、$\pi/(2\sqrt2 a)$、$\pi/(2a)$；留数与 30 位数值积分一致。}

\end{document}
